% part: intuitionistic-logic
% chapter: propositions-as-types
% section: reduction

\documentclass[../../../include/open-logic-section]{subfiles}

\begin{document}

\olfileid{int}{pty}{red}

\olsection{归约}

在自然演绎推导中，一个引入规则后紧跟着一个消去规则，就构成了一条可被“归约”的绕路。
例如，一个 $\Intro{\land}$ 后接 $\Elim{\land}$ 会“抵消”，因为 $\Elim{\land}$ 的结论总是 $\Intro{\land}$ 所得合取式的一个合取支，而 $\Intro{\land}$ 的两个前提即为这两个合取支。
因此，如果我们需要其中一个合取支的推导，可以直接使用那个推导，而不必先引入合取式再将其消去。我们有
\begin{gather*}
  \bottomAlignProof
  \AxiomC{}
  \RightLabel{$\delta_1$}
  \DeduceC{$!A_1$}
  \AxiomC{}
  \RightLabel{$\delta_2$}
  \DeduceC{$!A_2$}
  \RightLabel{$\Intro{\land}$}
  \BinaryInfC{$!A_1 \land !A_2$}
  \RightLabel{$\Elim{\land}_i$}
  \UnaryInfC{$!A_i$}
  \DisplayProof\bottomAlignProof
  \quad
  \redone
  \quad
  \AxiomC{}
  \RightLabel{$\delta_i$}
  \DeduceC{$!A_i$}
  \DisplayProof
\end{gather*}

由这种简化启发，我们可以在相应的证明项上考虑类似的归约关系。
如果 $N_1$ 是 $\delta_1$ 的证明项，而 $N_2$ 是 $\delta_2$ 的证明项，那么可以说，左侧推导的证明项（一步）归约到右侧推导的证明项：
\[
  \proj{i}{\pair{N_1, N_2}} \redone N_i
\]
在类型λ演算中，这就是积类型的 \emph{β-归约}规则。

在自然演绎中，像 \Intro{\land} 后接 \Elim{\land} 这样的一对推理（即可进行归约的一对推理）称为一个\emph{切割}。
在类型λ演算中，$\redone$ 左侧的项称为\emph{可约式}，右侧的项称为\emph{约化式}。
与无类型λ演算（其中只有 $(\lambd[x][N])Q$ 被视为可约式）不同，在类型λ演算中，语法扩展至包含 $\pair{N}{M}$、$\proj{i}{N}$、$\inj{i}{!A}{N}$、$\dcase{N}{x_1}{M_1}{x_2}{M_2}$ 和 $\abort{!A}{N}$ 的项，并具有相应的可约式。

类似地，我们有析取的归约：
\begin{gather*}
  \bottomAlignProof
  \AxiomC{}
  \RightLabel{$\delta$}
  \DeduceC{$!A_i$}
  \RightLabel{$\Intro{\lor}$}
  \UnaryInfC{$!A_1 \lor !A_2$}
  \AxiomC{$\Discharge{!A_1}{x_1}$}
  \RightLabel{$\delta_1$}
  \DeduceC{$!C$}
  \AxiomC{$\Discharge{!A_2}{x_2}$}
  \RightLabel{$\delta_2$}
  \DeduceC{$!C$}
  \DischargeRule{$\Elim{\lor}$}{x_1,x_2}
  \TrinaryInfC{$!C$}
  \DisplayProof\bottomAlignProof
  \quad
  \redone
  \quad
  \AxiomC{}
  \RightLabel{$\delta$}
  \DeduceC{$!A_i$}
  \RightLabel{$\delta_i$}
  \DeduceC{$!C$}
  \DisplayProof
\end{gather*}
这对应于证明项上的归约：
\begin{gather*}
  \dcase{\inj[!A_i]{i}{M}}{x_1}{N_1}{x_2}{N_2} \redone \Subst{N_i}{M}{x_i}
\end{gather*}
其中 $M$ 是 $\delta$ 的证明项，而 $N_i$ 是 $\delta_i$ 的证明项。这是和类型的 β-归约规则。
这里，$\Subst{M}{N_i}{x}$ 意味着将 $M$ 中变元~$x$ 的所有自由出现替换为~$N_i$。

函数类型的归约对应于消除 $\Intro\lif$ 后接 $\Elim\lif$ 的绕路。
\begin{gather*}
  \bottomAlignProof
  \AxiomC{$\Discharge{!A}{x}$}
  \RightLabel{$\delta$}
  \DeduceC{$!B$}
  \DischargeRule{$\Intro{\lif}$}{x}
  \UnaryInfC{$!A \lif !B$}
  \AxiomC{}
  \RightLabel{$\delta'$}
  \DeduceC{$!A$}
  \RightLabel{$\Elim{\lif}$}
  \BinaryInfC{$!B$}
  \DisplayProof\bottomAlignProof
  \quad
  \redone
  \quad
  \AxiomC{}
  \RightLabel{$\delta'$}
  \DeduceC{$!A$}
  \RightLabel{$\delta$}
  \DeduceC{$!B$}
  \DisplayProof
\end{gather*}
对于证明项而言，这等价于普通的 β-归约：
\begin{gather*}
  (\lambd[\typeof{x}{!A}][N]M)
  \redone \Subst{N}{M}{x}
\end{gather*}

除了推导上的归约变换，我们还可以考虑置换变换。
它们适用于以 \Elim{\lor} 后接另一消去规则结尾的（子）推导，例如：
\begin{multline*}
  \bottomAlignProof
  \AxiomC{}
  \RightLabel{$\delta$}
  \DeduceC{$!A_i$}
  \RightLabel{$\Intro{\lor}$}
  \UnaryInfC{$!A_1 \lor !A_2$}
  \AxiomC{$\Discharge{!A_1}{x_1}$}
  \RightLabel{$\delta_1$}
  \DeduceC{$!C_1 \land !C_2$}
  \AxiomC{$\Discharge{!A_2}{x_2}$}
  \RightLabel{$\delta_2$}
  \DeduceC{$!C_1 \land !C_2$}
  \DischargeRule{$\Elim{\lor}$}{x_1,x_2}
  \TrinaryInfC{$!C_1 \land !C_2$}
  \RightLabel{\Elim{\land}[i]}
  \UnaryInfC{$!C_i$}
  \DisplayProof\bottomAlignProof
  \quad
  \redone
  \\
  \AxiomC{}
  \RightLabel{$\delta$}
  \DeduceC{$!A_i$}
  \RightLabel{$\Intro{\lor}$}
  \UnaryInfC{$!A_1 \lor !A_2$}
  \AxiomC{$\Discharge{!A_1}{x_1}$}
  \RightLabel{$\delta_1$}
  \DeduceC{$!C_1 \land !C_2$}
  \RightLabel{\Elim{\land}[i]}
  \UnaryInfC{$!C_i$}
  \AxiomC{$\Discharge{!A_2}{x_2}$}
  \RightLabel{$\delta_2$}
  \DeduceC{$!C_1 \land !C_2$}
  \RightLabel{\Elim{\land}[i]}
  \UnaryInfC{$!C_i$}
  \DischargeRule{$\Elim{\lor}$}{x_1,x_2}
  \TrinaryInfC{$!C_i$}
  \DisplayProof
\end{multline*}
如果 $\delta$、$\delta_1$ 和 $\delta_2$ 的证明项分别是 $M$、$N_1$ 和 $N_2$，则证明项（即 $\lambd$ 项）相应的归约规则为：
\[
  \proj{i}{\dcase{M}{x_1}{N_1}{x_2}{N_2}} \redone
\dcase{M}{x_1}{\proj{i}{N_1}}{x_2}{\proj{i}{N_2}}
\]

当然，我们不仅可以消除推导末端的绕路，也可以将归约变换应用于以绕路结尾的子推导，从而消除推导内部的绕路。
类似地，$\beta$-归约可应用于 $\lambd$-项的子项。
如果 $N$ 是 $M$ 的一个子项，且 $N \redone N'$，我们可以用 $N'$ 替换 $M$ 中的子项 $N$ 以得到 $M'$。
在这种情况下，我们也记作 $M \redone M'$。
如果存在一个序列 $M = M_1
\redone M_2 \redone M_2 \redone \dots \redone M_k = M'$（包括 $k=1$ 即 $M = M'$ 的情形），我们记作 $M \red M'$。

无法对其任何子推导应用变换的推导称为\emph{正规的}。
类似地，无法（对其任何子项）应用变换的 $\lambd$-项也称为\emph{正规的}。
正规的 $\lambd$-项也称为\emph{范式}。

\end{document}